\begin{statementbox}
	\textbf{\textcolor{CStatement}{条件}}
	\begin{description}[style=nextline, leftmargin=2.8em, labelsep=0.5em, itemsep=0.35em]
		\item[\textbf{\textcolor{CStatement}{条件1（标准函数）：}}]
			函数形式为 $y=\sin x$, $y=\cos x$, $y=\tan x$，无其他参数。
	\end{description}
	\textbf{\textcolor{CStatement}{结论}}
	\begin{description}[style=nextline, leftmargin=2.8em, labelsep=0.5em, itemsep=0.45em]
		\item[\textbf{\textcolor{CStatement}{结论一（定义域）：}}]
			\textbf{\textcolor{CStatement}{直观描述：}}
			正弦和余弦函数的定义域为全体实数，正切函数的定义域为除去 $\frac{\pi}{2}+k\pi$ 的实数。
			\[
				D_{\sin}= \mathbb{R},\quad D_{\cos}= \mathbb{R},\quad D_{\tan}= \mathbb{R}\setminus\{\frac{\pi}{2}+k\pi\mid k\in\mathbb{Z}\}.
			\]
		\item[\textbf{\textcolor{CStatement}{结论二（值域）：}}]
			\textbf{\textcolor{CStatement}{直观描述：}}
			正弦和余弦函数的值域为 $[-1,1]$，正切函数的值域为 $\mathbb{R}$。
			\[
				R_{\sin}=[-1,1],\quad R_{\cos}=[-1,1],\quad R_{\tan}= \mathbb{R}.
			\]
		\item[\textbf{\textcolor{CStatement}{结论三（周期）：}}]
			\textbf{\textcolor{CStatement}{直观描述：}}
			正弦和余弦函数的最小正周期为 $2\pi$，正切函数的最小正周期为 $\pi$。
			\[
				T_{\sin}=2\pi,\quad T_{\cos}=2\pi,\quad T_{\tan}=\pi.
			\]
		\item[\textbf{\textcolor{CStatement}{结论四（奇偶性）：}}]
			\textbf{\textcolor{CStatement}{直观描述：}}
			正弦函数是奇函数，余弦函数是偶函数，正切函数是奇函数。
			\[
				\sin(-x)=-\sin x,\quad \cos(-x)=\cos x,\quad \tan(-x)=-\tan x.
			\]
		\item[\textbf{\textcolor{CStatement}{结论五（单调增区间）：}}]
			\textbf{\textcolor{CStatement}{直观描述：}}
			正弦、余弦、正切函数的单调递增区间如下：
			\[
				\begin{aligned}
					\sin &: \left[-\frac{\pi}{2}+2k\pi,\frac{\pi}{2}+2k\pi\right],\\\cos &: \left[-\pi+2k\pi,2k\pi\right],\\\tan &: \left(-\frac{\pi}{2}+k\pi,\frac{\pi}{2}+k\pi\right).
			\end{aligned}
			\]
	\end{description}
	\textbf{\textcolor{CStatement}{适用范围：}}
	以上性质仅适用于标准正弦、余弦、正切函数；对于复合函数需使用相应公式变形。
\end{statementbox}
